% Chapter 2 — API Obscape
\chapter{API Obscape}
\label{chap:api}

\section{Acceso}

\begin{table}[h!]
\centering
\begin{tabular}{ll}
\toprule
\textbf{Parámetro} & \textbf{Valor} \\
\midrule
Base URL         & \texttt{https://obscape.com/portal/api/v3/api} \\
Portal           & \texttt{https://obscape.com/portal} \\
Username         & \texttt{fuster} \\
Password portal  & \texttt{Delfos17*} \\
API Key          & \texttt{c1RyHhP6aJBPRHwIUrpz9eEPHPGhlbuMZIujEUvWTJaJPXJO0x} \\
\bottomrule
\end{tabular}
\caption{Credenciales Obscape}
\end{table}

\textbf{Nota:} password del portal y API key son cosas distintas. La API key
se obtiene en Portal $\to$ usuario $\to$ \emph{User Settings}.

El patrón base de cualquier llamada es:
\begin{Verbatim}
https://obscape.com/portal/api/v3/api?username=fuster&key=<API_KEY>&...
\end{Verbatim}

\section{Endpoints principales}

\begin{table}[h!]
\centering
\begin{tabular}{lp{8.5cm}}
\toprule
\textbf{Propósito} & \textbf{Parámetros adicionales} \\
\midrule
Listar proyectos                & (ninguno) \\
Listar estaciones del proyecto  & \texttt{\&project=<nombre>} \\
Datos últimas 24h               & \texttt{\&project=<p>\&station=<id>} \\
Datos entre fechas              & \texttt{\&station=<id>\&from=yyyy-mm-ddThh:mm:ss\&to=...} \\
Últimas N horas                 & \texttt{\&station=<id>\&latest=<N>} \\
Últimos N minutos               & \texttt{\&station=<id>\&latestMinutes=<N>} \\
Solo datos (sin metadatos)      & \texttt{\&dataonly} \\
Imagen por timestamp unix       & \texttt{\&station=<id>\&image=<unix\_ts>} \\
Última imagen disponible        & \texttt{\&station=<id>\&image=latest} \\
Timezone local                  & \texttt{\&tz=local} \\
\bottomrule
\end{tabular}
\caption{Endpoints de la API Obscape}
\end{table}

\subsection{Formato de nombre de imagen descargada}

\begin{Verbatim}
{unix_timestamp}_{YYYYMMDD}_{HHMMSS}_{station_id}.jpg
\end{Verbatim}
Ejemplo: \texttt{1777550400\_20260430\_120000\_PTM61474.jpg}

\section{Estado actual (verificado 2026-05-19)}

La API key funciona y devuelve:
\begin{lstlisting}[language=json]
[{
  "id":        "5866",
  "name":      "Ketel Haven",
  "devices":   [],
  "latitude":  "51.82553",
  "longitude": "4.727377",
  "reference": ""
}]
\end{lstlisting}

\textbf{Problema:} el proyecto registrado bajo la cuenta \texttt{fuster} es
``\emph{Ketel Haven}'', ubicado en los Países Bajos
(lat 51.82°N, lon 4.73°E), y tiene \texttt{"devices": []}. Las \textbf{6
cámaras de Guardamar del Segura no están vinculadas} a esta cuenta.
Todas las consultas de estación devuelven código \texttt{400}.

\textbf{Acción pendiente:} contactar con Obscape o el Ayuntamiento para
confirmar bajo qué cuenta/proyecto están registradas las cámaras.

\section{Script cliente --- \texttt{obscape\_api.py}}

Cliente Python en \texttt{obscape\_api.py} (raíz del proyecto). Uso:

\begin{Verbatim}
# Listar proyectos y estaciones
python obscape_api.py

# Descargar última imagen
python obscape_api.py --download

# Rango de fechas (solo 12:00h por defecto)
python obscape_api.py --from 2026-04-30 --to 2026-05-01

# Rango de fechas, todas las horas
python obscape_api.py --from 2026-04-30 --to 2026-05-01 --all-hours

# Últimas 24h de una estación concreta
python obscape_api.py --station PTM61474 --latest 24

# Directorio de salida personalizado
python obscape_api.py --download --out /ruta/salida
\end{Verbatim}

La clase \texttt{ObscapeClient} encapsula todos los métodos:
\texttt{list\_projects()}, \texttt{list\_stations()},
\texttt{get\_station\_data()}, \texttt{download\_image()},
\texttt{download\_range()}.

El método \texttt{download\_range()} filtra por hora
(\texttt{hour\_filter=12} por defecto) para priorizar las imágenes de las 12:00h.

\section{Dataset manual disponible}

Las imágenes de la estación \textbf{PTM61474} están disponibles en
\texttt{proces\_images/images/}: $\approx 90$ imágenes horarias
(04:00h--18:00h) entre el 29/04/2026 y el 06/05/2026.
Provienen del fichero ZIP \texttt{camera\_8214\_from20260429\_174800.zip}
facilitado manualmente.
